\section{Answers and explanations}
\label{sec:exercise-answers}

The numbers match the questions in \cref{sec:exercises}. Each answer identifies all correct
options and explains every choice. Question numbers link back to the exercises.

\Needspace{19\baselineskip}
\subsection{How an agent moves through a task}
\noindent Related material: \cref{sec:efsm}.

\begin{enumerate}[start=1,itemsep=0.7\baselineskip]
\item[{\hyperref[exercise:1]{1.}}] \begin{minipage}[t]{\linewidth}
\textbf{Correct: A, C.}
\begin{enumerate}[label=\textnormal{(\Alph*)},leftmargin=*,itemsep=0pt]
\item \textbf{Correct.} The step reads saved information, so a new observation can lead it down a different route.
\item \textbf{Incorrect.} Revisiting a step does not undo updates made by earlier steps.
\item \textbf{Correct.} The step name identifies the work to do; the saved values identify the particular situation.
\item \textbf{Incorrect.} The existing step can choose among its existing routes using changed information.
\end{enumerate}
\end{minipage}

\item[{\hyperref[exercise:2]{2.}}] \begin{minipage}[t]{\linewidth}
\textbf{Correct: B.}
\begin{enumerate}[label=\textnormal{(\Alph*)},leftmargin=*,itemsep=0pt]
\item \textbf{Incorrect.} The program's routing rule does not produce or control the external reply.
\item \textbf{Correct.} Different replies can satisfy different checks; each specific reply still has a definite route.
\item \textbf{Incorrect.} The difference can come entirely from different replies, without any random choice in the rule.
\item \textbf{Incorrect.} A routing rule selects an action after a reply arrives; it does not by itself model reply likelihood.
\end{enumerate}
\end{minipage}

\item[{\hyperref[exercise:3]{3.}}] \begin{minipage}[t]{\linewidth}
\textbf{Correct: D.}
\begin{enumerate}[label=\textnormal{(\Alph*)},leftmargin=*,itemsep=0pt]
\item \textbf{Incorrect.} Tool results can enter saved information and change later decisions.
\item \textbf{Incorrect.} Model replies can request different tools or request submission.
\item \textbf{Incorrect.} Allowed routes show possibilities, but they do not identify which one occurred.
\item \textbf{Correct.} With the same starting information and the same inputs at each step, the fixed rules select the same routes.
\end{enumerate}
\end{minipage}

\item[{\hyperref[exercise:4]{4.}}] \begin{minipage}[t]{\linewidth}
\textbf{Correct: A, B.}
\begin{enumerate}[label=\textnormal{(\Alph*)},leftmargin=*,itemsep=0pt]
\item \textbf{Correct.} The list changes the saved information, not the number of named steps.
\item \textbf{Correct.} Lists of arbitrarily many results can distinguish arbitrarily many situations.
\item \textbf{Incorrect.} The existing tool and decision steps can process another result.
\item \textbf{Incorrect.} An unbounded result list may make that table infinite.
\end{enumerate}
\end{minipage}

\item[{\hyperref[exercise:5]{5.}}] \begin{minipage}[t]{\linewidth}
\textbf{Correct: B, D.}
\begin{enumerate}[label=\textnormal{(\Alph*)},leftmargin=*,itemsep=0pt]
\item \textbf{Incorrect.} The same result occurs in both runs, but the call limit permits another call in only one of them.
\item \textbf{Correct.} The program needs that total to distinguish a run below the limit from one at the limit.
\item \textbf{Incorrect.} An accurate running total is enough for this limit; individual replies may be retained for other purposes.
\item \textbf{Correct.} The call count matters to the next admission check, while irrelevant details can be omitted from this decision.
\end{enumerate}
\end{minipage}

\item[{\hyperref[exercise:6]{6.}}] \begin{minipage}[t]{\linewidth}
\textbf{Correct: A, C, D.}
\begin{enumerate}[label=\textnormal{(\Alph*)},leftmargin=*,itemsep=0pt]
\item \textbf{Correct.} The program could not derive the particular reply from its saved information alone.
\item \textbf{Incorrect.} The check decides whether a route applies; the count changes through an update.
\item \textbf{Correct.} This step issues a request for tool execution. Later handling may check the request before the tool runs.
\item \textbf{Correct.} Those updates become part of the information retained between steps.
\end{enumerate}
\end{minipage}

\item[{\hyperref[exercise:7]{7.}}] \begin{minipage}[t]{\linewidth}
\textbf{Correct: C.}
\begin{enumerate}[label=\textnormal{(\Alph*)},leftmargin=*,itemsep=0pt]
\item \textbf{Incorrect.} A possible route is not evidence that it occurred.
\item \textbf{Incorrect.} The diagram assigns no probabilities, and path length does not supply one.
\item \textbf{Correct.} The diagram separates the route taken from other permitted continuations.
\item \textbf{Incorrect.} They remain possible for other inputs or runs, even though this run took the blue route.
\end{enumerate}
\end{minipage}

\item[{\hyperref[exercise:8]{8.}}] \begin{minipage}[t]{\linewidth}
\textbf{Correct: B, D.}
\begin{enumerate}[label=\textnormal{(\Alph*)},leftmargin=*,itemsep=0pt]
\item \textbf{Incorrect.} Many different texts have the same length.
\item \textbf{Correct.} Success is one of the facts the event explicitly records.
\item \textbf{Incorrect.} The log omits the result text and much of the agent's saved information.
\item \textbf{Correct.} The rules can explain which move connects recorded events, but that move was not logged directly.
\end{enumerate}
\end{minipage}

\item[{\hyperref[exercise:9]{9.}}] \begin{minipage}[t]{\linewidth}
\textbf{Correct: A, C.}
\begin{enumerate}[label=\textnormal{(\Alph*)},leftmargin=*,itemsep=0pt]
\item \textbf{Correct.} The plan is passed to the controller, whose replies can then change.
\item \textbf{Incorrect.} Subgoals are stored task data; the workflow's step names stay fixed.
\item \textbf{Correct.} Its input now includes a revised plan and possibly new observations.
\item \textbf{Incorrect.} In this agent, routing depends on replies and checks, not on which subgoal is current.
\end{enumerate}
\end{minipage}

\item[{\hyperref[exercise:10]{10.}}] \begin{minipage}[t]{\linewidth}
\textbf{Correct: B, D.}
\begin{enumerate}[label=\textnormal{(\Alph*)},leftmargin=*,itemsep=0pt]
\item \textbf{Incorrect.} The agent keeps the full history and checks only the results since revision.
\item \textbf{Correct.} After one revision, the stall route ends the run instead of revising again.
\item \textbf{Incorrect.} The workflow permits one revision; a later stall has a different outcome.
\item \textbf{Correct.} That position marks the start of the post-revision observation window.
\end{enumerate}
\end{minipage}

\end{enumerate}

\Needspace{19\baselineskip}
\subsection{What an agent remembers between steps}
\noindent Related material: \cref{sec:state}.

\begin{enumerate}[start=11,itemsep=0.7\baselineskip]
\item[{\hyperref[exercise:11]{11.}}] \begin{minipage}[t]{\linewidth}
\textbf{Correct: A, D.}
\begin{enumerate}[label=\textnormal{(\Alph*)},leftmargin=*,itemsep=0pt]
\item \textbf{Correct.} A step that omits a field leaves its stored value in place.
\item \textbf{Incorrect.} Steps return only fields they change; omission is not a deletion.
\item \textbf{Incorrect.} The runtime retains omitted fields without asking the node to repeat them.
\item \textbf{Correct.} The runtime applies the update to the named field and keeps the other stored values.
\end{enumerate}
\end{minipage}

\item[{\hyperref[exercise:12]{12.}}] \begin{minipage}[t]{\linewidth}
\textbf{Correct: D.}
\begin{enumerate}[label=\textnormal{(\Alph*)},leftmargin=*,itemsep=0pt]
\item \textbf{Incorrect.} Conversation length is not a reliable count, and the model does not enforce this limit.
\item \textbf{Incorrect.} The reply arrives after the call, so it cannot perform the pre-call check.
\item \textbf{Incorrect.} A history records events but does not itself reject a call.
\item \textbf{Correct.} The runtime can enforce the limit using saved information that is absent from the model's context.
\end{enumerate}
\end{minipage}

\item[{\hyperref[exercise:13]{13.}}] \begin{minipage}[t]{\linewidth}
\textbf{Correct: B, C.}
\begin{enumerate}[label=\textnormal{(\Alph*)},leftmargin=*,itemsep=0pt]
\item \textbf{Incorrect.} The log contains selected events rather than the full working information and pending work.
\item \textbf{Correct.} Workspace artifacts can survive even when the agent's in-memory information is lost.
\item \textbf{Correct.} It records selected actions and results, not a complete resumable state.
\item \textbf{Incorrect.} Files and in-memory information have different lifetimes.
\end{enumerate}
\end{minipage}

\item[{\hyperref[exercise:14]{14.}}] \begin{minipage}[t]{\linewidth}
\textbf{Correct: A, B, C.}
\begin{enumerate}[label=\textnormal{(\Alph*)},leftmargin=*,itemsep=0pt]
\item \textbf{Correct.} The repetition check needs earlier results, so replacing the list would discard needed evidence.
\item \textbf{Correct.} The new plan supersedes the old one for later work.
\item \textbf{Correct.} Each call contributes a delta to a total that lasts for the run.
\item \textbf{Incorrect.} That would lose the calls made before the latest one.
\end{enumerate}
\end{minipage}

\item[{\hyperref[exercise:15]{15.}}] \begin{minipage}[t]{\linewidth}
\textbf{Correct: A, C.}
\begin{enumerate}[label=\textnormal{(\Alph*)},leftmargin=*,itemsep=0pt]
\item \textbf{Correct.} An empty list contributes no new elements to an append operation.
\item \textbf{Incorrect.} Appending zero elements is different from replacing the stored list.
\item \textbf{Correct.} A reset must bypass the append rule and write the new list as the stored value.
\item \textbf{Incorrect.} An update changes a value; it does not remove the field declaration.
\end{enumerate}
\end{minipage}

\item[{\hyperref[exercise:16]{16.}}] \begin{minipage}[t]{\linewidth}
\textbf{Correct: B.}
\begin{enumerate}[label=\textnormal{(\Alph*)},leftmargin=*,itemsep=0pt]
\item \textbf{Incorrect.} A later route could treat the old reason as a new stall before the revised plan has been evaluated.
\item \textbf{Correct.} The signal has served its intended reader and should not influence a later routing decision.
\item \textbf{Incorrect.} The one-use signal should be cleared, but the tool-result history remains useful to later checks and validation.
\item \textbf{Incorrect.} Adding an empty list would leave the old reason in place, so it would not clear the active signal.
\end{enumerate}
\end{minipage}

\item[{\hyperref[exercise:17]{17.}}] \begin{minipage}[t]{\linewidth}
\textbf{Correct: A, B.}
\begin{enumerate}[label=\textnormal{(\Alph*)},leftmargin=*,itemsep=0pt]
\item \textbf{Correct.} That boundary separates results from before and after the revision.
\item \textbf{Correct.} Earlier repeated results should not count as new evidence of another stall.
\item \textbf{Incorrect.} They were produced under the previous plan and would cause an immediate false stall.
\item \textbf{Incorrect.} A saved boundary provides a fresh window while preserving history for other readers.
\end{enumerate}
\end{minipage}

\item[{\hyperref[exercise:18]{18.}}] \begin{minipage}[t]{\linewidth}
\textbf{Correct: A, C.}
\begin{enumerate}[label=\textnormal{(\Alph*)},leftmargin=*,itemsep=0pt]
\item \textbf{Correct.} The budget applies to the whole run, so earlier calls must still count.
\item \textbf{Incorrect.} This could admit more than 40 calls across several attempts.
\item \textbf{Correct.} Every place that can issue a call must honor the same limit.
\item \textbf{Incorrect.} A prompt does not supply a reliable program check before a call is issued.
\end{enumerate}
\end{minipage}

\item[{\hyperref[exercise:19]{19.}}] \begin{minipage}[t]{\linewidth}
\textbf{Correct: A, C.}
\begin{enumerate}[label=\textnormal{(\Alph*)},leftmargin=*,itemsep=0pt]
\item \textbf{Correct.} The stored count must match the calls actually completed for the check to be sound.
\item \textbf{Incorrect.} A call from another step could bypass the limit.
\item \textbf{Correct.} Changing it without a call would break its meaning as the completed-call total.
\item \textbf{Incorrect.} The reset would hide earlier calls and permit the run to exceed the whole-run limit.
\end{enumerate}
\end{minipage}

\item[{\hyperref[exercise:20]{20.}}] \begin{minipage}[t]{\linewidth}
\textbf{Correct: A, C.}
\begin{enumerate}[label=\textnormal{(\Alph*)},leftmargin=*,itemsep=0pt]
\item \textbf{Correct.} The declared field names give the static checker information to compare with the code.
\item \textbf{Incorrect.} The editor's static check does not install run-time validation for returned values.
\item \textbf{Correct.} A field marked optional in the declaration can still cause a missing-value problem if a running step expects it.
\item \textbf{Incorrect.} Partial updates leave existing values alone; they do not create values needed by an early reader.
\end{enumerate}
\end{minipage}

\end{enumerate}

\Needspace{19\baselineskip}
\subsection{Choosing the next action and knowing when to stop}
\noindent Related material: \cref{sec:control}.

\begin{enumerate}[start=21,itemsep=0.7\baselineskip]
\item[{\hyperref[exercise:21]{21.}}] \begin{minipage}[t]{\linewidth}
\textbf{Correct: B, C.}
\begin{enumerate}[label=\textnormal{(\Alph*)},leftmargin=*,itemsep=0pt]
\item \textbf{Incorrect.} The check receives saved information; seeing a result earlier does not make its discarded contents available later.
\item \textbf{Correct.} That is the only fact this step saved for the later check.
\item \textbf{Correct.} A check can use a relevant detail when an earlier step preserves it in the run's state; in this scenario, that detail was discarded.
\item \textbf{Incorrect.} Completion says the search ran; it does not say what the search found.
\end{enumerate}
\end{minipage}

\item[{\hyperref[exercise:22]{22.}}] \begin{minipage}[t]{\linewidth}
\textbf{Correct: B.}
\begin{enumerate}[label=\textnormal{(\Alph*)},leftmargin=*,itemsep=0pt]
\item \textbf{Incorrect.} A first-match rule selects one next step, even when several tests are true.
\item \textbf{Correct.} The program's test order resolves the overlap for this state.
\item \textbf{Incorrect.} The fallback is reached only when none of the earlier tests selects a step.
\item \textbf{Incorrect.} The priority is written into the program's test order, not inferred from the model's prose.
\end{enumerate}
\end{minipage}

\item[{\hyperref[exercise:23]{23.}}] \begin{minipage}[t]{\linewidth}
\textbf{Correct: A, D.}
\begin{enumerate}[label=\textnormal{(\Alph*)},leftmargin=*,itemsep=0pt]
\item \textbf{Correct.} The program examines whether the reply contains a request and uses that saved fact to choose a route.
\item \textbf{Incorrect.} The rule reads the request, not the persuasive wording of the reply.
\item \textbf{Incorrect.} A model reply supplies data to the rule; it does not rewrite the program's routing logic.
\item \textbf{Correct.} The rule sends an active reply containing a tool request to the tools step.
\end{enumerate}
\end{minipage}

\item[{\hyperref[exercise:24]{24.}}] \begin{minipage}[t]{\linewidth}
\textbf{Correct: B, C.}
\begin{enumerate}[label=\textnormal{(\Alph*)},leftmargin=*,itemsep=0pt]
\item \textbf{Incorrect.} A cycle can revisit a step within one compiled workflow.
\item \textbf{Correct.} The rule is fixed, while the values it reads can differ across runs.
\item \textbf{Correct.} A repeated location still has its own state and place in the run's sequence of decisions.
\item \textbf{Incorrect.} Compilation fixes possible routes; run-time values select one route at each decision.
\end{enumerate}
\end{minipage}

\item[{\hyperref[exercise:25]{25.}}] \begin{minipage}[t]{\linewidth}
\textbf{Correct: A, C, D.}
\begin{enumerate}[label=\textnormal{(\Alph*)},leftmargin=*,itemsep=0pt]
\item \textbf{Correct.} The next proposal can use that evidence to address the specific failure.
\item \textbf{Incorrect.} The repair path retains the conversation, including the failed proposal and reasons.
\item \textbf{Correct.} A repair needs a new candidate while preserving the evidence that explains the previous failure.
\item \textbf{Correct.} This repair path has no separate retry counter, so the run's call limit bounds repeated attempts.
\end{enumerate}
\end{minipage}

\item[{\hyperref[exercise:26]{26.}}] \begin{minipage}[t]{\linewidth}
\textbf{Correct: B, D.}
\begin{enumerate}[label=\textnormal{(\Alph*)},leftmargin=*,itemsep=0pt]
\item \textbf{Incorrect.} Without the recorded reason, the planner loses evidence of what failed and may repeat the same approach.
\item \textbf{Correct.} The planner can respond to the failure, and later checks then judge the new plan on its own evidence.
\item \textbf{Incorrect.} The workflow keeps observations; it changes the point from which the stall window is measured.
\item \textbf{Correct.} Only one revision is allowed, so a second stall ends the run instead of starting another revision.
\end{enumerate}
\end{minipage}

\item[{\hyperref[exercise:27]{27.}}] \begin{minipage}[t]{\linewidth}
\textbf{Correct: C.}
\begin{enumerate}[label=\textnormal{(\Alph*)},leftmargin=*,itemsep=0pt]
\item \textbf{Incorrect.} A fresh attempt must clear the old patch so it is not mistaken for the new candidate.
\item \textbf{Incorrect.} Plans are run data; replacing one does not change the workflow's structure.
\item \textbf{Correct.} Resetting either count would make used attempts or calls available again and could defeat their limits.
\item \textbf{Incorrect.} A reflection explains the failure but does not measure how many attempts have already been used.
\end{enumerate}
\end{minipage}

\item[{\hyperref[exercise:28]{28.}}] \begin{minipage}[t]{\linewidth}
\textbf{Correct: A, C.}
\begin{enumerate}[label=\textnormal{(\Alph*)},leftmargin=*,itemsep=0pt]
\item \textbf{Correct.} The model-call count does not change during that retry cycle.
\item \textbf{Incorrect.} A limit on the number of calls says nothing about the duration of one call or tool operation.
\item \textbf{Correct.} A decreasing count only bounds repeated cycles if each step between decreases finishes.
\item \textbf{Incorrect.} The check recognizes success when it occurs; it cannot force a future patch to pass.
\end{enumerate}
\end{minipage}

\item[{\hyperref[exercise:29]{29.}}] \begin{minipage}[t]{\linewidth}
\textbf{Correct: A, D.}
\begin{enumerate}[label=\textnormal{(\Alph*)},leftmargin=*,itemsep=0pt]
\item \textbf{Correct.} The observed failure pattern is irrelevant but varied, so it does not meet the identical-result condition.
\item \textbf{Incorrect.} Difference alone says nothing about relevance to the task.
\item \textbf{Incorrect.} A detector for one failure pattern supplies no success or termination guarantee for other patterns.
\item \textbf{Correct.} The budget applies to calls regardless of whether a stall pattern is detected.
\end{enumerate}
\end{minipage}

\item[{\hyperref[exercise:30]{30.}}] \begin{minipage}[t]{\linewidth}
\textbf{Correct: B, C, D.}
\begin{enumerate}[label=\textnormal{(\Alph*)},leftmargin=*,itemsep=0pt]
\item \textbf{Incorrect.} It shows that the acceptance procedure passed, which can still miss a defect in the patch.
\item \textbf{Correct.} Those are the two claims recorded by an accepting terminal event.
\item \textbf{Correct.} It could reflect an active run or a process stopped before recording a terminal event.
\item \textbf{Correct.} The terminal status and exception record carry different parts of the explanation.
\end{enumerate}
\end{minipage}

\end{enumerate}

\Needspace{19\baselineskip}
\subsection{Tool requests, execution, and access}
\noindent Related material: \cref{sec:boundary}.

\begin{enumerate}[start=31,itemsep=0.7\baselineskip]
\item[{\hyperref[exercise:31]{31.}}] \begin{minipage}[t]{\linewidth}
\textbf{Correct: A, C.}
\begin{enumerate}[label=\textnormal{(\Alph*)},leftmargin=*,itemsep=0pt]
\item \textbf{Correct.} Writes from this batch are applied after the workers finish, so finish order does not expose a partial update.
\item \textbf{Incorrect.} Updates remain invisible to other workers in the same batch until the batch's update phase.
\item \textbf{Correct.} The runtime applies the batch's writes before planning later work.
\item \textbf{Incorrect.} Competing writes need a specified way to combine them or must be rejected.
\end{enumerate}
\end{minipage}

\item[{\hyperref[exercise:32]{32.}}] \begin{minipage}[t]{\linewidth}
\textbf{Correct: B.}
\begin{enumerate}[label=\textnormal{(\Alph*)},leftmargin=*,itemsep=0pt]
\item \textbf{Incorrect.} The batch ordering rule governs saved state, not changes a tool already made to a file.
\item \textbf{Correct.} The external edit can happen during execution before the worker returns a state update.
\item \textbf{Incorrect.} A worker's writes become visible only when the runtime applies its completed update.
\item \textbf{Incorrect.} A file edit is an external effect and can occur before any update to saved run state is applied.
\end{enumerate}
\end{minipage}

\item[{\hyperref[exercise:33]{33.}}] \begin{minipage}[t]{\linewidth}
\textbf{Correct: A, D.}
\begin{enumerate}[label=\textnormal{(\Alph*)},leftmargin=*,itemsep=0pt]
\item \textbf{Correct.} Text shown to the model may influence its requests, but the execution path needs an access check to enforce the boundary.
\item \textbf{Incorrect.} A boundary must handle any request that reaches the tool, including a mistaken or hostile one.
\item \textbf{Incorrect.} Permission belongs to the runtime's policy, not to the model's account of why it wants a file.
\item \textbf{Correct.} Checking each file at the point of access closes the route that the description left open.
\end{enumerate}
\end{minipage}

\item[{\hyperref[exercise:34]{34.}}] \begin{minipage}[t]{\linewidth}
\textbf{Correct: C.}
\begin{enumerate}[label=\textnormal{(\Alph*)},leftmargin=*,itemsep=0pt]
\item \textbf{Incorrect.} A request can still supply the argument; an instruction does not remove that route.
\item \textbf{Incorrect.} A caller can override a default whenever the argument remains bindable.
\item \textbf{Correct.} Then the model has no open argument slot through which to replace the run setting.
\item \textbf{Incorrect.} A justification in the request does not transfer authority over trusted configuration.
\end{enumerate}
\end{minipage}

\item[{\hyperref[exercise:35]{35.}}] \begin{minipage}[t]{\linewidth}
\textbf{Correct: B, D.}
\begin{enumerate}[label=\textnormal{(\Alph*)},leftmargin=*,itemsep=0pt]
\item \textbf{Incorrect.} A correctly named path can still refer to a forbidden location.
\item \textbf{Correct.} Name binding and file permission are separate checks at different stages of execution.
\item \textbf{Incorrect.} Binding establishes that an argument fits the signature; the tool can still refuse it under the path policy.
\item \textbf{Correct.} Type correction handles a representation mistake; it does not change the access policy.
\end{enumerate}
\end{minipage}

\item[{\hyperref[exercise:36]{36.}}] \begin{minipage}[t]{\linewidth}
\textbf{Correct: A, C, D.}
\begin{enumerate}[label=\textnormal{(\Alph*)},leftmargin=*,itemsep=0pt]
\item \textbf{Correct.} The reason tells the model why the request failed and what to clarify.
\item \textbf{Incorrect.} The request did not identify which occurrence to change, so picking one risks an unintended edit.
\item \textbf{Correct.} A refusal must happen before the protected effect, leaving the file unchanged for that request.
\item \textbf{Correct.} A recoverable refusal lets the model refine its request and continue the run.
\end{enumerate}
\end{minipage}

\item[{\hyperref[exercise:37]{37.}}] \begin{minipage}[t]{\linewidth}
\textbf{Correct: B, C.}
\begin{enumerate}[label=\textnormal{(\Alph*)},leftmargin=*,itemsep=0pt]
\item \textbf{Incorrect.} A request can be refused or fail before the tool executes.
\item \textbf{Correct.} The result records success after dispatch rather than merely a proposed action.
\item \textbf{Correct.} A length reports how much text came back, not which characters it contained.
\item \textbf{Incorrect.} The request record precedes the checks; their outcome appears in the result.
\end{enumerate}
\end{minipage}

\item[{\hyperref[exercise:38]{38.}}] \begin{minipage}[t]{\linewidth}
\textbf{Correct: A, D.}
\begin{enumerate}[label=\textnormal{(\Alph*)},leftmargin=*,itemsep=0pt]
\item \textbf{Correct.} A tested check on one route leaves the new route able to bypass the boundary.
\item \textbf{Incorrect.} The new tool follows a different execution path and may skip the tested check.
\item \textbf{Incorrect.} The instruction may influence requests but does not enforce access at execution time.
\item \textbf{Correct.} Centralizing the check is effective when no alternate route reaches the file without it.
\end{enumerate}
\end{minipage}

\item[{\hyperref[exercise:39]{39.}}] \begin{minipage}[t]{\linewidth}
\textbf{Correct: B, C.}
\begin{enumerate}[label=\textnormal{(\Alph*)},leftmargin=*,itemsep=0pt]
\item \textbf{Incorrect.} That text also prefixes /tmp/ws-copy, and links or '..' can change the actual destination.
\item \textbf{Correct.} Resolution accounts for parent segments and symbolic links before deciding containment.
\item \textbf{Correct.} A link or parent segment can lead out even when the initial text appears local.
\item \textbf{Incorrect.} Joining an absolute path can leave it absolute; its resolved destination still needs the containment check.
\end{enumerate}
\end{minipage}

\item[{\hyperref[exercise:40]{40.}}] \begin{minipage}[t]{\linewidth}
\textbf{Correct: B, C, D.}
\begin{enumerate}[label=\textnormal{(\Alph*)},leftmargin=*,itemsep=0pt]
\item \textbf{Incorrect.} Search reads matching files; the editing tool is the writer and checks its target.
\item \textbf{Correct.} Search can return outside lines without checking each matched file against the workspace.
\item \textbf{Correct.} The check would refuse a match whose real location is outside the workspace before its contents are read.
\item \textbf{Correct.} The two tools reach files through different paths, so one passing test cannot establish the other's behavior.
\end{enumerate}
\end{minipage}

\end{enumerate}

\Needspace{19\baselineskip}
\subsection{Events and queued work}
\noindent Related material: \cref{sec:events}.

\begin{enumerate}[start=41,itemsep=0.7\baselineskip]
\item[{\hyperref[exercise:41]{41.}}] \begin{minipage}[t]{\linewidth}
\textbf{Correct: B, D.}
\begin{enumerate}[label=\textnormal{(\Alph*)},leftmargin=*,itemsep=0pt]
\item \textbf{Incorrect.} A queued request records intended work. The worker has not yet executed it.
\item \textbf{Correct.} Delivery to a worker registered for that kind of work is what starts the computation.
\item \textbf{Incorrect.} Putting an item in a queue does not perform the worker's file operations.
\item \textbf{Correct.} A record of execution or a successful result can show the checks ran; the queued request shows only intent. A refusal result would show that the checks did not run.
\end{enumerate}
\end{minipage}

\item[{\hyperref[exercise:42]{42.}}] \begin{minipage}[t]{\linewidth}
\textbf{Correct: A, C.}
\begin{enumerate}[label=\textnormal{(\Alph*)},leftmargin=*,itemsep=0pt]
\item \textbf{Correct.} Application code contains the steps that read the plan, update saved information, and choose what work comes next.
\item \textbf{Incorrect.} The function can put a new item in the queue; the runtime delivers it to the registered worker.
\item \textbf{Correct.} The loop checks pending items and calls the function registered for their kind.
\item \textbf{Incorrect.} Those rules come from the team's code, even though the runtime controls invocation.
\end{enumerate}
\end{minipage}

\item[{\hyperref[exercise:43]{43.}}] \begin{minipage}[t]{\linewidth}
\textbf{Correct: B, C.}
\begin{enumerate}[label=\textnormal{(\Alph*)},leftmargin=*,itemsep=0pt]
\item \textbf{Incorrect.} An empty queue only says that no work is scheduled; it does not record a successful outcome.
\item \textbf{Correct.} The completed plan was supposed to lead to another computation, but nothing can receive its message.
\item \textbf{Correct.} Work can cease because delivery failed, without a worker making a stopping decision.
\item \textbf{Incorrect.} Dispatch uses registered item kinds, and no registration exists for this one.
\end{enumerate}
\end{minipage}

\item[{\hyperref[exercise:44]{44.}}] \begin{minipage}[t]{\linewidth}
\textbf{Correct: C.}
\begin{enumerate}[label=\textnormal{(\Alph*)},leftmargin=*,itemsep=0pt]
\item \textbf{Incorrect.} A quiet queue can result from missing work as well as intentional completion.
\item \textbf{Incorrect.} The last worker could have failed to schedule its continuation.
\item \textbf{Correct.} A declared outcome is needed to distinguish success, a planned stop, and a lost continuation.
\item \textbf{Incorrect.} Queue state says nothing about the result of checks that may or may not have run.
\end{enumerate}
\end{minipage}

\item[{\hyperref[exercise:45]{45.}}] \begin{minipage}[t]{\linewidth}
\textbf{Correct: B, D.}
\begin{enumerate}[label=\textnormal{(\Alph*)},leftmargin=*,itemsep=0pt]
\item \textbf{Incorrect.} That message identifies the next kind of work but does not contain the full saved result.
\item \textbf{Correct.} The controller then reads the updated information when it starts.
\item \textbf{Incorrect.} The runtime may deliver the new item immediately, before the later save happens.
\item \textbf{Correct.} Repeating the same update could make one tool result appear twice in the saved history.
\end{enumerate}
\end{minipage}

\item[{\hyperref[exercise:46]{46.}}] \begin{minipage}[t]{\linewidth}
\textbf{Correct: C.}
\begin{enumerate}[label=\textnormal{(\Alph*)},leftmargin=*,itemsep=0pt]
\item \textbf{Incorrect.} Repeating a message is a delivery error and does not establish a second tool execution.
\item \textbf{Incorrect.} A repeated message does not undo an append to history.
\item \textbf{Correct.} Two deliveries can start duplicate handling of the same occurrence.
\item \textbf{Incorrect.} Duplicate delivery does not explain when or why the run should end.
\end{enumerate}
\end{minipage}

\item[{\hyperref[exercise:47]{47.}}] \begin{minipage}[t]{\linewidth}
\textbf{Correct: B, D.}
\begin{enumerate}[label=\textnormal{(\Alph*)},leftmargin=*,itemsep=0pt]
\item \textbf{Incorrect.} Silence provides no item for this loop to deliver.
\item \textbf{Correct.} That item gives the loop an event it can deliver to the timeout worker.
\item \textbf{Incorrect.} A log entry records an occurrence but is not necessarily an item submitted for dispatch.
\item \textbf{Correct.} Without a reply or timer event, no input arrives to start the timeout transition.
\end{enumerate}
\end{minipage}

\item[{\hyperref[exercise:48]{48.}}] \begin{minipage}[t]{\linewidth}
\textbf{Correct: D.}
\begin{enumerate}[label=\textnormal{(\Alph*)},leftmargin=*,itemsep=0pt]
\item \textbf{Incorrect.} The decision to end the run should prevent further work from starting.
\item \textbf{Incorrect.} An item requesting work does not determine the status of a run.
\item \textbf{Incorrect.} A declared stop is the point after which no new worker should start.
\item \textbf{Correct.} The terminal decision ends dispatch even if stray work remains.
\end{enumerate}
\end{minipage}

\item[{\hyperref[exercise:49]{49.}}] \begin{minipage}[t]{\linewidth}
\textbf{Correct: A, D.}
\begin{enumerate}[label=\textnormal{(\Alph*)},leftmargin=*,itemsep=0pt]
\item \textbf{Correct.} There is only one active continuation, so it maps to one point in the original workflow.
\item \textbf{Incorrect.} The workers may be at different parts of the workflow at the same time.
\item \textbf{Incorrect.} Both workers still read and update saved information, and their timing now matters more.
\item \textbf{Correct.} A single current-step label cannot represent simultaneous computations at different locations.
\end{enumerate}
\end{minipage}

\item[{\hyperref[exercise:50]{50.}}] \begin{minipage}[t]{\linewidth}
\textbf{Correct: A, C, D.}
\begin{enumerate}[label=\textnormal{(\Alph*)},leftmargin=*,itemsep=0pt]
\item \textbf{Correct.} These are the observable consequences of the transition rules at each comparison point.
\item \textbf{Incorrect.} The comparison fixes inputs and step boundaries, not scheduling overhead or timing.
\item \textbf{Correct.} Different messages can lead a live model to give different replies even when the step rules match.
\item \textbf{Correct.} The prepared comparison describes one continuation at a time, so concurrent cases require separate analysis.
\end{enumerate}
\end{minipage}

\end{enumerate}

\Needspace{19\baselineskip}
\subsection{Budgets and spending}
\noindent Related material: \cref{sec:budgets}.

\begin{enumerate}[start=51,itemsep=0.7\baselineskip]
\item[{\hyperref[exercise:51]{51.}}] \begin{minipage}[t]{\linewidth}
\textbf{Correct: A, C.}
\begin{enumerate}[label=\textnormal{(\Alph*)},leftmargin=*,itemsep=0pt]
\item \textbf{Correct.} A limit can prevent a new charge only if the permission check happens before work begins.
\item \textbf{Incorrect.} The bill does not exist yet, so it cannot be used for that call's starting decision.
\item \textbf{Correct.} Later starting decisions need an accurate total of charges already incurred.
\item \textbf{Incorrect.} The input grows with the conversation, and outputs can have different lengths.
\end{enumerate}
\end{minipage}

\item[{\hyperref[exercise:52]{52.}}] \begin{minipage}[t]{\linewidth}
\textbf{Correct: B, D.}
\begin{enumerate}[label=\textnormal{(\Alph*)},leftmargin=*,itemsep=0pt]
\item \textbf{Incorrect.} The 4 units already set aside must also be counted while the earlier call is unfinished.
\item \textbf{Correct.} Starting it would commit up to 11 units in total: 4 spent, 4 set aside, and 3 more.
\item \textbf{Incorrect.} Ignoring it would let multiple calls claim the same remaining capacity.
\item \textbf{Correct.} Releasing the unused 2 units would leave enough room to set aside 3 for a later call.
\end{enumerate}
\end{minipage}

\item[{\hyperref[exercise:53]{53.}}] \begin{minipage}[t]{\linewidth}
\textbf{Correct: B, C.}
\begin{enumerate}[label=\textnormal{(\Alph*)},leftmargin=*,itemsep=0pt]
\item \textbf{Incorrect.} The provider warns that actual input usage can differ from the estimate.
\item \textbf{Correct.} If the final charge can exceed the set-aside amount, the run can overspend when the call finishes.
\item \textbf{Correct.} Generation stops at the requested maximum, so output use cannot exceed that limit.
\item \textbf{Incorrect.} Past agreement can guide a forecast, but it does not establish a bound for every future request.
\end{enumerate}
\end{minipage}

\item[{\hyperref[exercise:54]{54.}}] \begin{minipage}[t]{\linewidth}
\textbf{Correct: C.}
\begin{enumerate}[label=\textnormal{(\Alph*)},leftmargin=*,itemsep=0pt]
\item \textbf{Incorrect.} Both may see the same available amount and each may decide that its call fits.
\item \textbf{Incorrect.} The spending limit could already have been exceeded by then.
\item \textbf{Correct.} The second worker then sees the first worker's commitment before deciding.
\item \textbf{Incorrect.} Equal limits do not make their combined possible cost fit the remaining budget.
\end{enumerate}
\end{minipage}

\item[{\hyperref[exercise:55]{55.}}] \begin{minipage}[t]{\linewidth}
\textbf{Correct: C.}
\begin{enumerate}[label=\textnormal{(\Alph*)},leftmargin=*,itemsep=0pt]
\item \textbf{Incorrect.} The provider may have processed part or all of the request before the timeout.
\item \textbf{Incorrect.} A zero bill is an unsupported guess while the actual charge is unknown.
\item \textbf{Correct.} The amount remains committed until the runtime can record or rule out the bill.
\item \textbf{Incorrect.} A failed attempt can still incur a charge, and earlier charges also remain part of the run.
\end{enumerate}
\end{minipage}

\item[{\hyperref[exercise:56]{56.}}] \begin{minipage}[t]{\linewidth}
\textbf{Correct: B, C.}
\begin{enumerate}[label=\textnormal{(\Alph*)},leftmargin=*,itemsep=0pt]
\item \textbf{Incorrect.} Output-only counting excludes the tokens sent into the model.
\item \textbf{Correct.} Input and output token counts both contribute to this unit.
\item \textbf{Correct.} Fresh input, cache reads, and output can have different prices before being combined into one charge.
\item \textbf{Incorrect.} One call can resend a much longer conversation than another.
\end{enumerate}
\end{minipage}

\item[{\hyperref[exercise:57]{57.}}] \begin{minipage}[t]{\linewidth}
\textbf{Correct: B, D.}
\begin{enumerate}[label=\textnormal{(\Alph*)},leftmargin=*,itemsep=0pt]
\item \textbf{Incorrect.} Those checks do not show that the replacement text is complete enough to apply safely.
\item \textbf{Correct.} The output-limit stop reason warns that a seemingly valid request may be cut off.
\item \textbf{Incorrect.} The model call still ran and consumed provider resources.
\item \textbf{Correct.} A retry is another paid call whose larger possible output needs its own permission check.
\end{enumerate}
\end{minipage}

\item[{\hyperref[exercise:58]{58.}}] \begin{minipage}[t]{\linewidth}
\textbf{Correct: A.}
\begin{enumerate}[label=\textnormal{(\Alph*)},leftmargin=*,itemsep=0pt]
\item \textbf{Correct.} Other roles then cannot spend the part reserved for the final decision.
\item \textbf{Incorrect.} The assessment is a model call and incurs usage too.
\item \textbf{Incorrect.} Earlier calls remain charged even after the patch reaches assessment.
\item \textbf{Incorrect.} A late call still needs permission under the same run limit.
\end{enumerate}
\end{minipage}

\item[{\hyperref[exercise:59]{59.}}] \begin{minipage}[t]{\linewidth}
\textbf{Correct: A, C.}
\begin{enumerate}[label=\textnormal{(\Alph*)},leftmargin=*,itemsep=0pt]
\item \textbf{Correct.} Clearing a failed candidate must not erase charges incurred while producing it.
\item \textbf{Incorrect.} Repeated attempts could then exceed the stated limit for the run.
\item \textbf{Correct.} Every model role can create a charge, so leaving one out understates run spending.
\item \textbf{Incorrect.} Failed calls and rejected attempts still consume resources.
\end{enumerate}
\end{minipage}

\item[{\hyperref[exercise:60]{60.}}] \begin{minipage}[t]{\linewidth}
\textbf{Correct: A, C, D.}
\begin{enumerate}[label=\textnormal{(\Alph*)},leftmargin=*,itemsep=0pt]
\item \textbf{Correct.} Holding these conditions fixed helps isolate the effect of how work is scheduled.
\item \textbf{Incorrect.} A seed does not make different model inputs produce the same replies, and one trial cannot show variation.
\item \textbf{Correct.} Keeping tools inside each fresh workspace blocks access to files left by earlier trials. Outcomes and costs show how and why runs stopped.
\item \textbf{Correct.} Another model or budget could change both the trajectory and the result.
\end{enumerate}
\end{minipage}

\end{enumerate}

